\chapter{Conclusions}
\label{ch:conclusions}

% Allow a little extra line stretch so long technical compounds rebreak
% cleanly instead of spilling into the margin (no template packages added).
\setlength{\emergencystretch}{3em}

This dissertation set out to automate the most repetitive and least-supported phase
of a security audit, the revalidation of reported findings, without
surrendering the human judgement that a compliance-grade verdict demands. The result,
\emph{revalid}, is a deliberately narrow tool: it does not discover or exploit
vulnerabilities, but re-executes the reproduction steps of an already-reported finding
under a per-command human gate and returns a conservative, evidence-linked verdict.
Whether that specialisation earns its keep is ultimately an empirical question. The
first three objectives are realised in the artefact: penetration-test reports are
analysed into findings and their reproduction steps, those steps are interpreted into
gated verification commands, and a human-supervised retest turns them into an
evidence-backed verdict. The fourth, a quantitative evaluation of verdict
reliability and of the effort saved over manual retesting, is addressed by the study of
Chapter~\ref{ch:evaluation}, which exercises the whole pipeline against a real report on a
live target: eight of twelve findings correctly confirmed on a small local model, four
hedged, no false clearance on any of them, and a machine time of the same order as a
manual pass. A wider study across more reports, more targets and repeated runs per
finding is the direction those figures point; what the cross-model comparison already
shows is that the safety of the result rests on the approval gate rather than on the
capability of the model behind it. Those runs exercised the gate programmatically
(Section~\ref{sec:eval-threats}), so validating it with human operators is the first item
of the wider study. This closing chapter steps back from the artefact
to set down what the work \emph{taught}, and how the degree made it possible.

\section{Lessons learned}
\label{sec:lessons}

The largest single thing I take from this project is a working command of
\emph{AI-assisted software development}: not merely prompting a model, but
engineering an environment in which a capable, fast-moving assistant can be directed
reliably. Over the course of the work, Claude~Code moved from a code generator to the
medium through which the entire lifecycle was managed: requirements captured as
tracked issues, significant decisions recorded as architecture decision records,
recurring procedures encoded as project-specific \emph{skills}, quality checks
delegated to review \emph{agents}, and governance enforced automatically by
\emph{hooks} in the editing loop. Learning to build and tune that apparatus, and,
just as importantly, learning to recognise when it had grown heavier than the task
justified and to cut it back (Section~\ref{subsec:process-evolution}), is a
competence I did not have at the outset. These tools are moving quickly from novelty
to professional expectation, and having driven one end to end, across every layer of
a real project, is the learning I expect to carry furthest.

A second lesson concerned discipline over autonomy. The very reliability failures that
shape \emph{revalid}'s design, hallucination and over-confident conclusions, apply
equally to the assistant that builds it. What made AI assistance trustworthy here was
ordinary engineering rigour: a strict typing and test gate that no generated code
could bypass, a continuous-integration barrier before every merge, and a human review
on every change. The assistant accelerated the work; it never got to decide whether
the work was correct. I also came away with a far deeper understanding of the domain
the tool serves: the vulnerability-management lifecycle, the safety pitfalls of
language-model agents, and the genuinely hard problem of turning a human-written
report into something a machine can act on.

\section{Future work}
\label{sec:future}

Several directions remain open, from the workflow this tool would slot into, through
the two capability limits the evaluation exposed, to lifting the single-operator
scope the work deliberately adopted.

One of them comes first, because the central result rests on it. The evaluation showed
that a verdict authority separate from the retesting agent refuses the false clearance an
autonomous agent produces --- but it exercised that authority programmatically, and with
the expected verdicts known
(Section~\ref{sec:eval-threats}), whereas the design's claim is that a \emph{person}
should hold it. Testing that is a human-factors study rather than a software one:
operators who do not know the answer, working from the agent's evidence and its
recommendation, measured on how often they accept a confident but unsupported
\emph{fixed} --- and, in passing, on what the gate costs in a person's time, which this
evaluation could only bracket. Until it is run, the per-command human gate is a justified
design choice with its mechanism validated and its human half unmeasured.

\subsection{Closing the loop with automated remediation}
\label{subsec:future-remediation}

The tightest opportunity is to close the loop between remediation and revalidation.
As automated code-fixing agents such as CodeMender~\cite{codemender} begin to
\emph{apply} fixes, an automated, human-gated tool to \emph{confirm} those fixes
becomes a natural complement rather than a curiosity: the fix and its verification
are the two halves of the same workflow, and only one of them is being automated
today. A revalidation tool that already refuses to declare anything fixed without
evidence is the right shape for that second half.

\subsection{Verifying browser-observable findings}
\label{subsec:future-browser}

A sharper limitation surfaced in the evaluation of
Chapter~\ref{ch:evaluation}: a whole class of vulnerabilities can be confirmed only by
observing its effect take shape in a rendered page, which a purely command-line, HTTP
retest cannot see. Stored cross-site scripting and server-side template injection are
the clearest cases, with DOM-based and reflected scripting alongside them: a retest
over HTTP can establish that a payload is stored or reflected, but not that the injected
script \emph{executes} in the browser or that the template expression is
\emph{evaluated} on render. Faced with that gap the system did what NFR-01 demands:
it returned an \emph{inconclusive} verdict rather than guessing \emph{fixed}, a safe
hedge noted among the threats to validity (Section~\ref{sec:eval-threats}) rather than a
confident call.

The natural improvement is to give the retest sandbox eyes. A headless browser, driven
inside the same network-isolated sandbox and under the same per-command human gate as
every other action, would let the agent load the response, run its scripts and watch
whether the payload actually fires, turning a principled hedge into a confident
confirmation on exactly the findings the current design must otherwise leave open.
Browser-driven verification in fact existed in an earlier iteration of the project and
was retired when the batch architecture it belonged to was removed; reviving it as one
more tool the agent can reach for, now inside the agentic sandbox, would recover that
capability on the disciplined foundation the rest of the system was rebuilt on.

\subsection{Observing the model's own work}
\label{subsec:future-observability}

A third direction follows from a limitation this work states rather than solves, and
the evaluation of Chapter~\ref{ch:evaluation} is what makes it pressing. Reproducibility
(NFR-02) is met at the level of the \emph{verdict}: the session transcript records every
command, its real output and the determination drawn from them, and the audit re-derives
each stored verdict from it. What the trail does \emph{not} capture is the model's side
of each exchange: the exact prompt, the parameters, the token usage and the latency of
every individual call. That gap is tolerable for the claim actually made, since a verdict
is re-derivable without it, but it limits precisely the questions the study
had to answer by hand or leave open: why a particular turn proposed the command it did,
whether the four hedged findings failed on reasoning or on context length, and what a
retest costs in tokens and time on one backend against another, the comparison
Section~\ref{sec:eval-time} could only make in wall-clock, and only by naming the model
loaded at the time.

Closing that gap does not require building anything, which is what makes it attractive.
Pydantic~AI already emits OpenTelemetry spans for every agent step, so an
LLM-observability platform such as Langfuse, open-source and
self-hostable and therefore compatible with the localhost-only posture of NFR-03 and
with the privacy argument of Section~\ref{subsec:eval-local}, can ingest them with an
instrumentation call rather than a redesign. A wider evaluation instrumented that way
would report not only how often the system is right, but what each verdict cost and
which prompt produced it, turning the backend comparison from an observation into a
measurement. It is worth being clear about what such a platform would and would not add.
It is a development and analysis instrument, not a second audit trail: the append-only
transcript remains the record a verdict is derived from, because the reproducibility
argument of Section~\ref{sec:design-data} depends on that record being written by the
system itself rather than by an optional external collector.

\subsection{Growing the application}
\label{sec:scaling}

The tool was deliberately scoped as a local, single-operator application, a
decision that traded reach for the focus a thesis needs
(Section~\ref{subsec:prior-attempts}). The most natural line of future development is
to lift that restriction and grow it from a single-operator tool into a multi-user,
client--server platform, this time on the disciplined foundation the project
established, rather than the over-broad one it set out from.

Several upgrades follow from that goal. The single process that today serves one
operator would separate into a hosted backend and thin clients, with authentication
and role-based access control so that auditors, developers and reviewers can share an
engagement with appropriately different permissions. The single-file SQLite store
would move to a full relational database, a migration the data-access layer was
deliberately designed to absorb without a rewrite (Section~\ref{subsec:stack}), and
the agent's sandboxed executions would be dispatched to a pool of workers through a
job queue, so that many findings, or many engagements, can be revalidated concurrently
rather than one at a time. On top of that would sit the team-oriented features a
production tool needs: shared dashboards of open and closed findings, a cross-user
audit trail, and integration with the issue trackers and continuous-integration
pipelines where remediation actually happens, so that a revalidation can be triggered
automatically the moment a fix is deployed. Each of these is an increment on the
existing architecture rather than a rewrite, which is exactly the property the
disciplined, single-operator core was built to preserve.

\medskip
Beyond the artefact itself, what I value most from this work is the confidence that I
can take a technology as new and as powerful as an agentic coding assistant and still
hold it to the engineering standards the degree instilled, to use it fully but on
my own terms. That, more than the tool, is what I carry forward.
